\scanpage{55}
et $x$ est défini modulo $f$ indépendamment de $s$ et on trouve ainsi une section canonique $x_0$ de $L/fL$. En fait, \deleted{$L/fL=$} sur $S'_0=\unsure{\mathrm{Var}(f)|\mathrm{Van}(f)}=V(f\underline O_S)=V(2\underline O_S)$, on trouve ici que $x_0=\underline b_{0}(\Delta|S'_0)\in\Gamma(S'_0,L|S'_0)$ (d'après § 9). Ainsi, le torseur $F$ est déterminé \inserted{(si $2\,1_S$ est $\mathcal O_S$-régulier)} par la seule connaissance de \deleted{\unsure{$B$}} $L$ et de la section
\[b_0\in\Gamma(S_0,L_0)\tag*{(105)}\]

De façon précise, la donnée d'un torseur $F$ sous $\underline L$, et d'une section $\underline b$ de $fF$, équivaut \unsure{exactement|exactemt} ($f$ étant une section $\underline O_S$-régulière de $\underline O_S$) à la donnée d'une section $b_0$ de $L_0$, en lui associant le \inserted{(plutôt)} sous-faisceau Zariski de $\frac1f L_{\unsure{0}}$ formé des sections \inserted{locales} $x$ de $\frac1f L$ telles que
% slanted
\margin{marche pour $L$ faisceau quasi-coh. tel que $f$ soit $L$-réguli\unsure{er|ère}\ldots}
\[fx|S_0=b_0.\]

\deleted{$(F,\underline b)$ étant ainsi déterminés par la connaissance de $(L,b_0)$, \inserted{donc $B$ aussi} il reste à déterminer les $q$.}

Dans le cas actuel, on devra avoir
\[b_0^2=D_0\qquad({:}\,D|S_0)\tag*{(106)}\]
ce qui, si $S_0$ est réduit, détermine $b_0$ de façon unique [et l'existence de $b_0$ est une
% archive sheet number: 54
